\title{Política de Comunicação

do

Conselho Regional de

Psicologia de São Paulo}
\author{}
\date{}

\begin{document}
\maketitle

\chapter{Apresentação}

A comunicação institucional é indispensável para a autarquia e
desenvolve o papel de manter uma comunicação de duas vias, que não só
transmite informações como também promove a interação entre as pessoas e
dá suporte à governança na tomada de decisões, além de desenvolver e
consolidar relacionamentos.

Essas políticas contribuem para as boas práticas de confiabilidade da
governança, estabelecendo canais de comunicação efetivos que estimulam a
interação e facilitam a interlocução com os seus públicos.

A \textbf{Política de Comunicação} do CRP~SP é resultado do
\href{https://transparencia.cfp.org.br/crp06/planejamento/planejamento-estrategico-2024/}{planejamento
estratégico de 2024}} e cumpre o previsto na
\href{https://transparencia.cfp.org.br/crp06/legislacao/regimento-interno-do-conselho-regional-de-psicologia-de-sao-paulo-6a-regiao-aprovado-pela-resolucao-no-5-de-22-de-marco-de-2023/}{Resolução
CRP~SP nº~05/2023}}, que dispõe sobre o regimento da autarquia.

Acredita-se que, baseados nesses aspectos, os processos de comunicação
permitirão estabelecermos um canal direto entre trabalhadores, gestores,
conselheiras, colaboradoras e categoria profissional que permita o fluxo
de informação em todos os sentidos.

As ações de comunicação institucional a serem conduzidas serão reguladas
e orientadas por esta \textbf{Política}, que relaciona as atividades a
serem realizadas pela sede e subsedes.

Esta \textbf{Política}, elaborada pela Coordenadoria de Comunicação,
deverá ser publicada amplamente.

De acordo com esse direcionamento, todos são responsáveis pela
comunicação correta e eficaz, não sendo esta atribuição exclusiva da
Unidade de Comunicação.

É responsabilidade de todos que atuam no CRP~SP zelar pela boa imagem da
instituição e cuidar para que os processos de comunicação institucional
se realizem adequadamente, conforme os objetivos institucionais.

Portanto, é importante a sensibilização de todas as equipes,
independentemente da área em que atuam.

A presente \textbf{Política} tem por finalidade ser um instrumento
orientador e normativo de ações alinhadas aos planos estratégicos. Desta
maneira, a \textbf{Política de Comunicação} constitui um valioso
instrumento a ser observado e praticado por todos, independentemente do
vínculo com o CRP~SP.

Conselho Regional de Psicologia de São Paulo

XVII Plenária

\chapter{Introdução}

Nenhuma norma é capaz de prever todas as situações em que será
empregada. Quando se trata de uma área tão dinâmica e tão dependente da
tecnologia como a comunicação, a previsibilidade é ainda menor: novas
mídias, linguagens e plataformas surgem a todo o momento. Por isso,
nesta introdução apresentaremos alguns princípios gerais que devem ser
aplicados tanto a casos não previstos nesta \textbf{Política} quanto na
elaboração de novas versões deste documento.

\section{Uma política de comunicação
pública}

No Brasil, a ideia de comunicação pública ganha relevância a partir da
redemocratização, quando passa a substituir o conceito de ``comunicação
governamental'', introduzido em nossa administração pública durante os
governos de Getúlio Vargas (1930--1945). O paradigma de comunicação
pública é, assim, uma tentativa de superar a comunicação unilateral,
característica de governos autoritários, por meio da qual a população é
informada, mas não participa.

Para
\href{https://abcpublica.org.br/wp-content/uploads/2021/02/Comunica\%C3\%A7\%C3\%A3o-P\%C3\%BAblica-VF-Cap\%C3\%ADtulo.pdf}{Jorge
Duarte}}, ``fazer comunicação pública é assumir a perspectiva cidadã na
comunicação envolvendo temas de interesse coletivo, alterando seu eixo,
tradicionalmente centrado no atendimento dos interesses da instituição e
de seus gestores''. O objetivo da comunicação pública, conclui o autor,
é ``o atendimento do interesse público e da sociedade, simbolizado pelo
cidadão''.

De maneira mais didática, a Associação Brasileira de Comunicação Pública
(ABCPública) elenca
\href{https://abcpublica.org.br/conheca-os-12-principios-da-comunicacao-publica/}{doze
princípios da comunicação pública}}:

\begin{enumerate}

\item
  garantir o acesso amplo à informação;
\item
  fomentar o diálogo;
\item
  estimular a participação;
\item
  promover os direitos e a democracia;
\item
  combater a desinformação;
\item
  ouvir a sociedade;
\item
  focar no cidadão;
\item
  ser inclusiva e plural;
\item
  tratar a comunicação como política de Estado;
\item
  garantir a impessoalidade;
\item
  pautar-se pela ética;
\item
  atuar com eficácia.
\end{enumerate}

\section{Princípios
institucionais}

A comunicação do CRP~SP está pautada pelos princípios da ética, da
responsabilidade pública, da transparência, da credibilidade e da
colaboração. Materializa a voz e \emph{persona} institucionais e, por
meio dela, a autarquia mantém diálogo constante com a categoria e a
sociedade. Empenha-se em comunicar a diversidade e a representatividade
(racial, étnica, de gênero, de corpos, etária, de territórios). Seus
produtos adotam linguagem inclusiva (simples, gendrada, antirracista,
anticapacitista e com recursos de acessibilidade).

Os princípios institucionais a seguir estão presentes no
\href{https://correio.crpsp.org.br/service/home/~/?auth=co&loc=pt_BR&id=4201&part=6}{Regimento
Interno}} do CRP~SP e são derivados das normativas referentes ao Sistema
Conselhos de Psicologia e, principalmente, do
\href{https://www.crpsp.org/uploads/pagina/289379/2j9LIMPLJ9jFr5YrK57HmAiBWjVMxdbe.pdf}{Código
de ética profissional da/o psicóloga/o}}. Além destes, publicações de
caráter jornalístico estão sujeitas ao
\href{https://fenaj.org.br/codigo-de-etica-dos-jornalistas-brasileiros/}{Código
de ética dos jornalistas brasileiros}}.

Também são documentos relacionados à \textbf{Política de Comunicação} e
devem ser observados no desenvolvimento das atividades:

\begin{enumerate}

\item
  \href{https://www.planalto.gov.br/ccivil_03/constituicao/constituicao.htm}{Constituição
  Federal}};
\item
  \href{https://www.planalto.gov.br/ccivil_03/leis/l5766.htm}{Lei
  nº~5.766/71}}, que criou o Sistema Conselhos de Psicologia;
\item
  Regimento Interno do CRP~SP
  (\href{https://transparencia.cfp.org.br/crp06/legislacao/regimento-interno-do-conselho-regional-de-psicologia-de-sao-paulo-6a-regiao-aprovado-pela-resolucao-no-5-de-22-de-marco-de-2023/}{Resolução
  CRP~SP nº~05/2023}});
\item
  Código de Ética Profissional da/o Psicóloga/o
  (\href{https://atosoficiais.com.br/cfp/resolucao-do-exercicio-profissional-n-10-2005-aprova-o-codigo-de-etica-profissional-do-psicologo}{Resolução
  CFP n º~010/2005}});
\item
  \href{https://www.planalto.gov.br/ccivil_03/decreto/d70274.htm}{Decreto
  nº~70.274, de 9 de março de 1972}} (normas do cerimonial público e a
  ordem geral de precedência);
\item
  Lei de Acesso à Informação (LAI;
  \href{https://www.planalto.gov.br/ccivil_03/_ato2011-2014/2011/lei/l12527.htm}{Lei
  nº~12.527/2011}});
\item
  \href{https://dados.gov.br/dados/conteudo/politica-de-dados-abertos}{Política
  de Dados Abertos}} do Governo Federal;
\item
  \href{https://transparencia.cfp.org.br/wp-content/uploads/sites/29/2024/08/Politica-de-TI-2-Edicao.pdf}{Política
  de Segurança em Tecnologia da Informação e Comunicação}} do CRP~SP;
\item
  \href{https://transparencia.cfp.org.br/crp06/legislacao/portaria-crp-n-49-2024-institui-a-politica-de-patrocinio-e-apoio-institucional-do-conselho-regional-de-psicologia-da-6a-regiao-crp-06/}{Política
  de Patrocínio e Apoio Institucional}} do Conselho Regional de
  Psicologia da 6ª Região;
\item
  \href{http://cedoc.crpsp.org.br/handle/1/2938}{Manual de Linguagem
  do CRP~SP}}.
\end{enumerate}

\subsection{Respeito e promoção dos Direitos
Humanos}

A centralidade da promoção dos Direitos Humanos e do respeito à
dignidade da pessoa está expressa no Código de Ética Profissional da/o
Psicóloga/o e no artigo 130 do
\href{https://correio.crpsp.org.br/service/home/~/?auth=co&loc=pt_BR&id=4201&part=6}{Regimento
Interno}} do CRP~SP.

Além de prezar pela divulgação de pautas relacionadas ao tema, a
comunicação da autarquia deve demonstrar seu compromisso com os Direitos
Humanos em todas as suas práticas. A \textbf{linguagem} adotada pelo
Conselho, portanto, deve ser acima de tudo \textbf{inclusiva}: deve
respeitar as orientações da norma culta sem deixar de dar visibilidade a
grupos social e historicamente marginalizados.

Assim, entendemos por linguagem inclusiva a linguagem \textbf{gendrada}
(isto é, sem o uso do gênero masculino para generalizações),
\textbf{antirracista} (que não se utilize de palavras, expressões ou
figuras de linguagem de conotação preconceituosa ou baseadas em
estereótipos de raça ou etnia) e \textbf{antietarista} (que não se
utilize de palavras, expressões ou figuras de linguagem que tenham como
premissa uma visão negativa do envelhecimento).

A linguagem que adotamos também é \textbf{anticapacitista} --- não se
utiliza de palavras, expressões ou figuras de linguagem que atribuam
valor a competências cognitivas, sensoriais ou motoras. Tão importante
quanto usarmos uma linguagem que inclua as pessoas com deficiência,
contudo, é dispormos de \textbf{recursos de acessibilidade}: textos que
sejam lidos adequadamente por leitores de tela, \emph{sites} de
navegação intuitiva, descrição de imagens e audiodescrição e uso de
\textbf{linguagem simples}.

Nossos produtos são publicados em diversos formatos: no \emph{site},
como livro digital, na forma de postagens. Cada um desses formatos tem
diretrizes específicas de acessibilidade. \textbf{Em qualquer caso,
deverão ser seguidas as orientações do
``\href{https://emag.governoeletronico.gov.br/}{Modelo de
Acessibilidade em Governo Eletrônico}}`` (eMAG)}.

Como princípios fundamentais e complementares da comunicação do CRP~SP,
\textbf{inclusão} e \textbf{acessibilidade} devem guiar toda e qualquer
produção da unidade de comunicação e todo e qualquer projeto de
comunicação integrada.

\subsection{Transparência}

O artigo 131 do
\href{https://correio.crpsp.org.br/service/home/~/?auth=co&loc=pt_BR&id=4201&part=6}{Regimento
Interno}} do CRP~SP prevê a divulgação dos atos do Conselho com o
objetivo de aumentar o conhecimento e o reconhecimento da Psicologia
pela sociedade brasileira.

\subsection{Orientação}

O artigo 132 do
\href{https://correio.crpsp.org.br/service/home/~/?auth=co&loc=pt_BR&id=4201&part=6}{Regimento
Interno}} do CRP~SP determina que a autarquia mantenha ``publicações
destinadas à divulgação de matéria de interesse da psicóloga e do
público em geral, cabendo ao Plenário e/ou Diretoria dispor a respeito e
de acordo com a Política de Comunicação Institucional''.

\subsection{Efetivação de funções precípuas do Conselho Regional de
Psicologia}

São funções precípuas dos conselhos profissionais, conforme o
\href{about:blank}{Tribunal de Contas da União}} (TCU), cabendo à
Coordenação de Comunicação colaborar com seu desempenho:

\begin{enumerate}

\item
  orientar, supervisionar, fiscalizar e disciplinar o exercício da
  profissão;
\item
  zelar pelo prestígio e bom nome da profissão e dos profissionais que a
  exercem;
\item
  organizar e manter o registro profissional (registro, inscrição,
  cancelamento);
\item
  examinar reclamações e representações;
\item
  representar às autoridades competentes;
\item
  atuar como órgão de consulta do governo;
\item
  atuar em colaboração com entidades de classe, escolas e faculdades;
\item
  contribuir para o desenvolvimento econômico;
\item
  disseminar a técnica econômica nos diversos setores da economia
  nacional, promovendo estudos e campanhas em prol da racionalização
  econômica do país;
\item
  promover estudos e campanhas em prol da racionalização administrativa
  do país.
\end{enumerate}

\section{Princípios constitucionais aplicáveis à comunicação
pública}

O
\href{https://www.planalto.gov.br/ccivil_03/decreto-lei/del0200.htm}{Decreto-Lei
nº~200}}, de 25 de fevereiro de 1967, define \textbf{autarquia} como

\begin{quote}
{[}\ldots{]} o serviço autônomo, criado por lei, com personalidade
jurídica, patrimônio e receita próprios, para executar atividades
típicas da Administração Pública, que requeiram para seu melhor
funcionamento, gestão administrativa e financeira descentralizada.
\end{quote}

Como autarquia, portanto, o CRP~SP está subordinado às normas e
princípios que regem a administração pública.

Por sua vez, o artigo 37 da
\href{https://www.planalto.gov.br/ccivil_03/constituicao/constituicao.htm}{Constituição
Federal}} prevê que ``A administração pública direta e indireta de
qualquer dos Poderes da União, dos Estados, do Distrito Federal e dos
Municípios obedecerá os princípios de legalidade, impessoalidade,
moralidade, publicidade e eficiência''. Esses princípios, recepcionados
no artigo 3º, inciso XXI do
\href{https://correio.crpsp.org.br/service/home/~/?auth=co&loc=pt_BR&id=4201&part=6}{Regimento
Interno do CRP~SP}}, são elucidados a seguir.

\subsubsection{Legalidade}

Diferentemente de instituições privadas, às quais é permitido fazer tudo
aquilo que não seja expressamente proibido em lei, às autarquias é
permitido fazer \textbf{apenas} o que a lei prevê.

A obediência ao princípio da legalidade deve, ainda, levar em conta a
hierarquia do ordenamento jurídico brasileiro, que tem a Constituição
Federal como lei maior, e abaixo dela as leis ordinárias --- das quais
destacamos a
\href{https://www.google.com/url?sa=t&source=web&rct=j&opi=89978449&url=https://www.planalto.gov.br/ccivil_03/leis/l5766.htm&ved=2ahUKEwi7voyM_MKJAxWYF7kGHcbCO54QFnoECBMQAw&usg=AOvVaw3BsZ1ugMxVxz5tlw9ZAIyv}{Lei
nº~5.766/1971}}, de 20 de dezembro de 1971, que cria o Conselho Federal
e os Conselhos Regionais de Psicologia, e a
\href{https://www.google.com/url?sa=t&source=web&rct=j&opi=89978449&url=https://www.planalto.gov.br/ccivil_03/leis/l9649cons.htm&ved=2ahUKEwjLibuZ_MKJAxVuILkGHaPnFtcQFnoECAoQAQ&usg=AOvVaw1dF6QYD5NfkI-TBa6TTJgo}{Lei
nº~9.649/1998}}, que dispõe, entre outras coisas, sobre os serviços de
fiscalização de profissões regulamentadas --- e as normativas próprias
da autarquia (portarias, resoluções etc.).

\subsubsection{Impessoalidade}

O princípio da impessoalidade determina que a atuação dos órgãos da
administração pública tem por finalidade atender ao interesse público.
No caso da comunicação do CRP~SP, isso implica tanto o atendimento das
demandas de profissionais de Psicologia quanto a defesa da categoria
perante a sociedade. A comunicação institucional da autarquia deve,
portanto, ser útil às psicólogas e aos psicólogos e responsável em sua
relação com a população em geral, pois representa a face oficial da
Psicologia no estado de São Paulo.

Como consequência do compromisso assumido com o interesse público, é
vedada a utilização da \emph{persona} institucional para a promoção de
indivíduos ou grupos. A comunicação institucional não pode fugir ao
princípio da legalidade, devendo restringir-se às situações previstas
nas normativas da instituição e na legislação pátria; segundo o § 1º~do
artigo 37 da
\href{https://www.google.com/url?sa=t&source=web&rct=j&opi=89978449&url=https://www.planalto.gov.br/ccivil_03/constituicao/constituicao.htm&ved=2ahUKEwjlis2B_MKJAxXcHrkGHTtTAvcQFnoECBgQAQ&usg=AOvVaw3i_8717crw9PBIV4q9Jndm}{Constituição
Federal}},

\begin{quote}
A publicidade dos atos, programas, obras, serviços e campanhas dos
órgãos públicos deverá ter caráter educativo, informativo ou de
orientação social, dela não podendo constar nomes, símbolos ou imagens
que caracterizem promoção pessoal de autoridades ou servidores públicos.
\end{quote}

\subsubsection{Moralidade}

O conceito administrativo de moralidade está ligado às ideias de
probidade, decoro e boa-fé. Nesse sentido, ``moralidade'' não é um
conceito subjetivo, mas um referencial de atuação adequada aos
pressupostos éticos adotados por um dado grupo social. Na administração
pública, esses pressupostos são definidos pela legislação que define os
limites de atuação de uma agente pública.

\subsubsection{Publicidade}

Como princípio da administração pública, a publicidade se refere tanto à
exigência de publicação de atos administrativos para que produzam
efeitos externos --- a publicação de uma portaria, por exemplo ---
quanto à exigência de transparência da atuação administrativa. É
necessário, por exemplo, que os atos administrativos apresentem sua
motivação, isto é, que explicitem os pressupostos fáticos e jurídicos
que o determinam, para que sociedade e órgãos de controle possam
fiscalizar a legitimidade desses atos.

\subsubsection{Eficiência}

Por fim, o princípio de eficiência fala tanto da qualidade da atuação da
agente pública quanto da racionalidade da administração. A eficiência
das agentes é determinada pela qualidade e produtividade de seu
trabalho; a racionalidade da administração está ligada à prestação de
serviço público adequado, definido pela
\href{https://www.planalto.gov.br/ccivil_03/leis/l8987cons.htm}{Lei
nº~8.987/95} como aquele que ``satisfaz as condições de regularidade,
continuidade, eficiência, segurança, atualidade, generalidade, cortesia
na sua prestação e modicidade das tarifas'' (art. 6º, § 1º).

\section{Implementação e gestão desta
Política}

Esta \textbf{Política} deverá ser observada por conselheiras e
conselheiros, bem como por trabalhadoras e trabalhadores do CRP~SP.

Questões não previstas nesta \textbf{Política} deverão ser encaminhadas
para o \emph{e-mail}
\href{mailto:ascom@crpsp.org.br}{ascom@crpsp.org.br}, para
avaliação e proposição às instâncias competentes.

\chapter{Estrutura e organização}

De acordo com a
\href{https://transparencia.cfp.org.br/crp06/wp-content/uploads/sites/29/2022/09/Resolucao-03-22-PECS-f-1-assinada.pdf}{Resolução
CRP nº~03, de 29 de agosto de 2022}, a estrutura administrativa do
Conselho possui duas instâncias decisórias: gerências e coordenações.

A Coordenação de Comunicação está subordinada à Gerência de Relações
Institucionais. É responsável por elaborar e executar o planejamento de
comunicação definido periodicamente com a Comissão de Comunicação
Institucional (ComCom).

\section{Comissão de Comunicação Institucional
(ComCom)}

A Comissão de Comunicação Institucional (ComCom) é uma comissão
permanente, que funciona como instância deliberativa e consultiva para
os assuntos relativos à comunicação do CRP~SP.

\subsection{Recursos}

A ComCom é composta de pelo menos duas conselheiras, tendo uma delas
como coordenadora.

\subsection{Atribuições}

De acordo com o
\href{https://correio.crpsp.org.br/service/home/~/?auth=co&loc=pt_BR&id=4201&part=6}{Regimento
Interno do CRP~SP}}, compete à ComCom:

\begin{enumerate}

\item
  acompanhar e aprovar as ações da Unidade de Comunicação Institucional
  do CRP-06, informando o Plenário;
\item
  responsabilizar-se pela comunicação institucional interna e externa;
\item
  responsabilizar-se pela produção e divulgação dos eventos
  institucionais externos de âmbito estadual;
\item
  dar diretrizes para a produção e divulgação dos eventos do CRP-06;
\item
  apresentar ao Plenário a identidade visual e o plano de comunicação
  das campanhas, ações e unidades do CRP-06;
\item
  elaborar a Política de Comunicação Institucional para aprovação do
  Plenário e fazer cumprir suas previsões;
\item
  acompanhar o adequado preenchimento e publicação do Portal da
  Transparência de acordo com o disposto na Lei de Acesso à Informação
  (LAI; Lei nº~12.527/2011) e na Resolução CFP nº~20/2018, ou outras que
  vierem a substituí-las;
\item
  propor pautas e acompanhar a elaboração do Jornal Psi para aprovação
  do Plenário;
\item
  criar o plano anual de comunicação institucional do CRP-06 em
  observância às priorizações do planejamento estratégico;
\item
  acompanhar e alimentar o site e redes sociais do CRP-06, encaminhando
  propostas e/ou denúncias, que ferem a imagem da autarquia, para
  providências da Diretoria, quando necessário;
\item
  realizar reuniões semanais com suas/seus membras/os para planejar,
  encaminhar e avaliar as ações;
\item
  articular com a Comissão de Comunicação do Conselho Federal de
  Psicologia;
\item
  supervisionar a edição das publicações do CRP-06;
\item
  coordenar a divulgação das ações do CRP-06;
\item
  acolher e encaminhar à unidade responsável para que responda às
  solicitações e consultas de pessoas físicas e jurídicas em relação às
  informações de domínio comum da entidade no site e redes sociais, bem
  como aquelas correlatas que lhe sejam atribuídas pelo Plenário;
\item
  organizar e manter atualizados os canais de comunicação do CRP-06;
\item
  orientar o diálogo com a imprensa e outros canais externos para
  participação do CRP-06 em suas pautas;
\item
  garantir a aplicação da linguagem gendrada, antirracista e inclusiva
  no CRP-06, em seus canais de comunicação internos e externos;
\item
  acompanhar e apresentar ao Plenário os resultados quantitativos e
  qualitativos da Comissão de Comunicação;
\item
  aprovar custos envolvendo fornecedoras/es e aquisição de outros
  recursos para a plena realização do planejamento e melhoria da área,
  acompanhando as licitações demandadas pela área;
\item
  propor, avaliar e encaminhar para aprovação do Plenário, sempre que
  necessário, produtos de mídia que possam ser usados para comunicação
  com as profissionais de Psicologia.
\end{enumerate}

\section{Gerência de Relações
Institucionais}

Tem a finalidade de gerir, avaliar e acompanhar a execução das
atividades das coordenações subordinadas, colaborando com a diretoria e
exercendo as competências que lhe forem delegadas.

\subsection{Recursos}

A Gerência de Relações Institucionais é ocupada por uma trabalhadora em
cargo de comissão.

\subsection{Atribuições}

Compete à Gerência de Relações Institucionais, segundo a
\href{https://transparencia.cfp.org.br/crp06/wp-content/uploads/sites/29/2022/09/Resolucao-03-22-PECS-f-1-assinada.pdf}{Resolução
CRP nº~03, de 29 de agosto de 2022}}:

\begin{enumerate}

\item
  executar as atividades relacionadas à comunicação e articulação com os
  segmentos da sociedade envolvidos com os serviços fiscalizados pelo
  CRP-06;
\item
  articular e desenvolver o relacionamento do CRP-06 com órgãos de
  interesse, tais como universidades, organismos nacionais e
  internacionais e organizações da sociedade civil;
\item
  executar atividades de divulgação de conteúdos, participação e
  realização de eventos relacionados ao trabalho desempenhado pelo
  CRP-06 e de interesse a suas/seus inscritas/os, divulgar conteúdos de
  interesse do CRP-06;
\item
  mapear, articular e acompanhar possíveis parcerias e convênios com
  órgãos relevantes aos fins do CRP-06, submetendo as propostas à
  Diretoria;
\item
  atuar de forma articulada com as demais Gerências responsabilizando-se
  pelas deliberações relativas aos processos e à gestão do trabalho no
  cumprimento das metas institucionais;
\item
  organizar, articular, acompanhar, avaliar e monitorar as atividades
  das unidades sob sua responsabilidade, utilizando-se de indicadores,
  quando necessário, para verificação das necessidades e garantia do bom
  desenvolvimento dos trabalhos; e
\item
  desempenhar outras atribuições correlatas de nível similar.
\end{enumerate}

\section{Coordenação de
Comunicação}

A Coordenação de Comunicação é a instância executiva da comunicação do
CRP~SP.

\subsection{Recursos}

A Coordenação de Comunicação é composta de:

\begin{enumerate}

\item
  uma coordenadora;
\item
  duas analistas em gestão de comunicação;
\item
  quatro profissionais de suporte administrativo (PSA);
\item
  duas estagiárias.
\end{enumerate}

\subsection{Atribuições}

De acordo com a
\href{https://transparencia.cfp.org.br/crp06/wp-content/uploads/sites/29/2022/09/Resolucao-03-22-PECS-f-1-assinada.pdf}{Resolução
CRP nº~03, de 29 de agosto de 2022}}, compete à Coordenação de
Comunicação:

\begin{enumerate}

\item
  planejar, organizar, coordenar e divulgar as ações referentes à
  integração das estratégias institucionais com a categoria e a
  sociedade;
\item
  construir, desenvolver e implantar estratégias com foco na aproximação
  com a categoria e a sociedade;
\item
  garantir a divulgação das atividades e posicionamentos institucionais
  do CRP-06 para a categoria e a sociedade;
\item
  gerenciar os servidores de \emph{e-mail} e hospedagem de \emph{sites}
  do CRP-06;
\item
  criar, manter e atualizar as aplicações \emph{web};
\item
  desenvolver trabalhos de artes gráficas que sejam demandados pelo
  CRP-06, acompanhando as produções e a aprovação do material
  apresentado pelos fornecedores (prova impressa, qualidade de impressão
  e material produzido);
\item
  fazer a cobertura de eventos com a utilização de filmagens,
  fotografias ou quaisquer outros meios necessários, assim como promover
  a sonorização, projeção de mídia e edição subsequentes necessárias;
\item
  produzir e diagramar publicações do CRP-06 (impressas ou eletrônicas),
  para divulgação às/aos psicólogas/os, contemplando assuntos de seus
  interesses, quando for o caso;
\item
  assessorar e desenvolver políticas de relações públicas,
  relacionando-se com as áreas de formação de opinião pública e
  diferentes mídias;
\item
  desenvolver e executar políticas de assessoria de imprensa,
  acompanhando o alcance e impacto das estratégias de divulgação,
  comunicação e diálogo com a categoria e a sociedade;
\item
  propor estratégia e divulgar ações técnico-políticas e posicionamentos
  do Conselho visando qualificar e publicizar o CRP-06;
\item
  divulgar interna e externamente as atividades desenvolvidas pelas
  unidades organizacionais do CRP-06;
\item
  coordenar e acompanhar a realização de cobertura jornalística
  técnico-política do CRP-06;
\item
  desenvolver e manter atualizado o banco de imagens e fontes do CRP-06;
\item
  criar e atualizar \emph{mailing} de contatos de pessoas físicas e
  entidades do CRP-06;
\item
  realizar todos os tipos de eventos (seminários, congressos, fóruns,
  debates e simpósios, entre outros) do CRP-06, contemplando todas as
  atividades necessárias que percorrem as etapas de planejamento,
  organização, realização e pós-evento;
\item
  fazer a cobertura de eventos com a utilização de filmagens,
  fotografias ou quaisquer outros meios necessários, assim como promover
  a sonorização, projeção de mídia e edição subsequentes necessárias;
\item
  manter e organizar os equipamentos audiovisuais do auditório do
  CRP-06;
\item
  apoiar a realização de eventos indicados pelo CRP-06 para outras
  entidades;
\item
  desenvolver e manter atualizado o banco de imagens e fontes do CRP-06;
\item
  realizar exposições do CRP-06 em espaços externos ao CRP-06 em toda a
  sua jurisdição;
\item
  operacionalizar concursos culturais do CRP-06 no estado de São Paulo;
\item
  operacionalizar o deslocamento e hospedagem, se necessário, de
  coordenadoras/es, palestrantes e convidadas/os para eventos do CRP-06;
\item
  pesquisar fornecedores e elaborar cotações para locação de serviços
  e/ou compra de materiais para eventos do CRP-06;
\item
  pesquisar e apresentar, para avaliação das instâncias superiores,
  espaços para a realização de eventos, analisando as adequações de
  localização e infraestrutura;
\item
  operar a mesa de som da cabine do auditório do CRP-06, incluindo
  microfones, DVD \emph{players} e demais equipamentos, realizando,
  inclusive, gravação e edição de áudio dos eventos;
\item
  atender as solicitações diversas das/os psicólogas/os, respondendo por
  \emph{e-mail}, por telefone e pessoalmente, sobre informações de
  eventos do CRP-06;
\item
  oferecer apoio administrativo na participação do CRP-06 em eventos
  promovidos por outras organizações, fazendo cotações de espaço,
  materiais, equipamentos e participando da promoção de publicações e/ou
  serviços em feiras e outros eventos no estado de São Paulo, prestando
  orientação às/aos visitantes e/ou participantes;
\item
  criar subsídios para a elaboração de documentos para apoio e
  patrocínio para a realização de eventos realizados pelo CRP-06;
\item
  identificar e criar estratégias para utilização de novos meios de
  comunicação;
\item
  desempenhar outras atribuições correlatas de nível similar.
\end{enumerate}

\subsubsection{Atribuições funcionais}

Compete à
\href{https://transparencia.cfp.org.br/crp06/wp-content/uploads/sites/29/2022/09/Resolucao-03-22-PECS-Anexo-II-Regulamentacao-Plano-de-Empregos-Carreiras-e-Salarios-f-assinada.pdf}{\textbf{coordenadora}}}
a gestão de processos e equipes de trabalho da Coordenação de
Comunicação do CRP~SP.

Compete às
\href{https://transparencia.cfp.org.br/crp06/wp-content/uploads/sites/29/2022/09/Resolucao-03-22-PECS-Anexo-II-Regulamentacao-Plano-de-Empregos-Carreiras-e-Salarios-f-assinada.pdf}{\textbf{analistas
em gestão de comunicação}}}:

\begin{enumerate}

\item
  selecionar, revisar, preparar e distribui materiais para publicação;
\item
  fotografar e gravar imagens jornalísticas;
\item
  editar publicações impressas e eletrônicas;
\item
  selecionar, divulgar e arquivar a comunicação feita a respeito do
  Conselho nos meios impressos e eletrônicos;
\item
  manter contato com a imprensa externa fornecendo dados, materiais e
  marcando entrevistas;
\item
  implantar ações de relações públicas e assessoria de imprensa;
\item
  organizar eventos internos e externos;
\item
  preparar, organizar, coordenar e realizar o cerimonial; e
\item
  desempenhar outras atribuições correlatas de nível similar.
\end{enumerate}

Compete às
\href{https://transparencia.cfp.org.br/crp06/wp-content/uploads/sites/29/2023/03/Resolucao-CRP-002-2023.pdf}{\textbf{profissionais
de suporte administrativo}}} \textbf{(\emph{design} gráfico)}:

\begin{enumerate}

\item
  preparar, organizar, elaborar, atualizar e conferir documentos,
  relatórios, gráficos e tabelas bem como prestar atendimento ao
  profissional e ao público em geral na área em que atua;
\item
  executar, sob supervisão, atividades de preparo técnico de comunicação
  interna e externa, gráficas utilizando a suíte de aplicativos de
  aplicativos Adobe (InDesign, Illustrator, Photoshop), CorelDraw e
  aplicativos similares;
\item
  preparar composições gráficas, diagramação, desenvolver marcas,
  materiais impressos, papelaria, sinalização, interface de web,
  infográficos e ilustrações;
\item
  definir especificações técnicas para produção de materiais impressos;
\item
  acompanhar a produção dos materiais analisando e aprovando provas
  gráfica e conteúdo;
\item
  fazer diagramação de livros, jornais e publicações similares;
\item
  realizar outras atividades correlatadas da unidade de lotação.
\end{enumerate}

Compete às
\href{https://transparencia.cfp.org.br/crp06/wp-content/uploads/sites/29/2023/03/Resolucao-CRP-002-2023.pdf}{\textbf{profissionais
de suporte administrativo}}}
\href{https://transparencia.cfp.org.br/crp06/wp-content/uploads/sites/29/2022/09/Resolucao-03-22-PECS-Anexo-II-Regulamentacao-Plano-de-Empregos-Carreiras-e-Salarios-f-assinada.pdf}{\textbf{(assistente
administrativo)}}:

\begin{enumerate}

\item
  executar trabalhos de recepção e suporte administrativo que envolvam:
  serviços de informação e controle de visitantes e correspondências,
  redação de documentos e relatórios, digitação, reprodução xerográfica,
  coleta, expedição, distribuição e arquivamento de documentos;
\item
  realizar visitas de expedientes junto a fóruns e comarcas;
\item
  requisitar, receber e controlar o uso de material de consumo;
\item
  realizar atividades de atendimento ao profissional, bem como
  atividades de caráter administrativo que envolva redação e digitação;
\item
  preparar e dar suporte administrativo às sessões plenárias éticas e de
  julgamento;
\item
  atender às solicitações das comissões técnicas ou grupos de trabalho
  específicos criados pelo CRP-06;
\item
  efetuar o registro das solicitações, interpretar e encaminhar ao
  demandante, como subsídios para as ações do CRP-06;
\item
  manter sistema de informação atualizado de todas as solicitações
  recebidas e encaminhadas para conhecimento interno e da direção do
  CRP-06;
\item
  verificar a situação profissional dos psicólogos contratados por
  instituições ou empresas que prestem serviços de psicologia;
\item
  fazer atendimento presencial e por telefone aos usuários dos serviços
  do CRP~SP;
\item
  preparar, organizar, atualizar, classificar, arquivar, elaborar,
  redigir, conferir e expedir documentos que envolva processos de
  registro, transferência, cancelamento da categoria de Psicólogos;
\item
  informar processos e fornecer subsídios para análise e tomada de
  decisão da diretoria;
\item
  controlar e manter atualizado o arquivamento dos prontuários de
  Psicólogos;
\item
  dar suporte administrativo às seções plenárias ordinárias, diretoria,
  comissões temáticas, grupos de trabalho e núcleos;
\item
  providenciar o deslocamento e hospedagem dos conselheiros de plenário,
  diretoria, grupos de trabalho, comissões e núcleos;
\item
  preparar, organizar, atualizar, elaborar e conferir documentos,
  relatórios, minutas de atos administrativos e normativos, gráficos e
  tabelas;
\item
  executar rotinas, sistemas, métodos e procedimentos de trabalho
  relativos ao processamento de folha de pagamentos, recolhimento de
  encargos sociais, rotinas de administração de pessoal, recrutamento e
  seleção e administração de benefícios;
\item
  executar processos de admissão, rescisão e administração de contratos
  de trabalho;
\item
  executar rotinas, sistemas, métodos e procedimentos de trabalho
  relativos a administração de fornecedores, contratos públicos,
  materiais, processo de compras, recebimento e expedição de materiais,
  controle de estoque, análise de custos, preços e resultados, análise
  de viabilidade econômica, administração e controle orçamentário;
\item
  preparar, organizar, atualizar, elaborar e conferir documentos,
  relatórios, gráficos e tabelas bem como prestar atendimento ao
  profissional e ao público em geral;
\item
  controlar fundo fixo (pequeno caixa), verbas, contas a pagar e fluxo
  de caixa, conferindo notas fiscais/recibos e prestando contas;
\item
  realizar conciliação e análise de movimentos contábeis;
\item
  fazer lançamentos contábeis da folha de pagamentos, balancetes e
  outros demonstrativos de contas preparando / elaborando a declaração
  de imposto de renda, DIRF e demonstrativos gerenciais, cadastramento e
  controle do arquivo inativo e do inventário físico;
\item
  realizar outras atividades correlatas da unidade de lotação.
\end{enumerate}

\chapter{Canais, serviços e
produtos}

\section{Canais oficiais}

Os canais oficiais de representação do Conselho Regional de Psicologia
de São Paulo (CRP~SP) são os seguintes:

\begin{enumerate}

\item
  \emph{site} e \emph{e-mail} institucional;
\item
  perfis oficiais em redes sociais;
\item
  eventos promovidos pela autarquia, em parceria ou não, nos estados e
  territórios.
\end{enumerate}

\subsection{Site institucional}

O \emph{site} institucional é o principal repositório de informações do
CRP~SP. Nele devem ser disponibilizadas informações sobre todas as
atividades do Conselho, além de material de referência e orientação para
a categoria. O \emph{site} também deve oferecer canais para contato com
todas as unidades da autarquia.

\subsection{Redes sociais}

O CRP~SP mantém perfis ativos nas seguintes redes sociais:

\begin{enumerate}

\item
  Facebook;
\item
  Instagram;
\item
  LinkedIn;
\item
  X/Twitter;
\item
  YouTube.
\end{enumerate}

Os perfis sociais do CRP~SP têm como objetivo formar e informar a
categoria sobre temas associados à Psicologia, à atuação profissional e
ao funcionamento do Sistema Conselhos. Por meio deles, buscamos ampliar
nossos pontos de contato e comunicação com psicólogas de todo o estado
de São Paulo, além de participar dos debates sobre a Psicologia como
ciência e profissão.

Em observância aos princípios da transparência e a nosso compromisso
ético, estabelecemos as seguintes diretrizes:

\begin{enumerate}

\item
  \textbf{perfis e páginas oficiais}: só são permitidos perfis e páginas
  associados ao CRP~SP que sejam criados e mantidos pela instituição. A
  autarquia se reserva o direito de não endossar perfis e páginas que se
  autointitulem associados ao Conselho e que não sejam reconhecidos como
  oficiais, tomando as devidas medidas quando necessário;
\item
  \textbf{gestão das contas}: as contas oficiais do CRP~SP serão
  gerenciadas por uma equipe certificada e autorizada, composta por
  profissionais de comunicação, trabalhadoras e conselheiras da
  autarquia;
\item
  \textbf{política de interação e moderação}: para garantir a
  veracidade, exatidão, confiabilidade, qualidade e legalidade das
  informações, as interações nos canais oficiais estarão sujeitas a
  moderação. Mensagens e perfis poderão ser excluídos se:

  \begin{enumerate}
  
  \item
    usarem linguagem inapropriada, obscena, caluniosa, grosseira,
    abusiva, difamatória ou ofensiva;
  \item
    fizerem apologia a práticas ilícitas;
  \item
    incitarem ódio, violência, racismo ou qualquer forma de
    discriminação;
  \item
    contiverem ameaças, assédio, injúria, calúnia, difamação ou qualquer
    outro ilícito penal;
  \item
    divulgarem conteúdos na forma de \emph{spam} ou ``correntes'';
  \item
    tiverem objetivo comercial ou publicitário;
  \item
    forem ininteligíveis ou fora de contexto;
  \item
    contiverem propagandas político-partidárias;
  \item
    contiverem \emph{links} suspeitos ou representarem ameaça à
    segurança da informação;
  \item
    usarem informações e imagens de pessoas e instituições de modo
    indevido;
  \item
    contiverem dados pessoais da autora, do autor ou de terceiros;
  \item
    violarem direitos de imagem e de propriedade intelectual;
  \item
    forem fraudulentos (perfis falsos e/ou \emph{bots}) ou promoverem
    conteúdo inverídico (\emph{fake news}).
  \end{enumerate}
\item
  \textbf{restrições de divulgação}: o CRP~SP não segue e não marca
  pessoas, veículos de comunicação e empresas privadas: apenas entidades
  parceiras e órgãos do poder público. Não divulgamos partidos políticos
  e não realizamos atendimento administrativo e de orientação nas redes
  sociais. As demandas que não se enquadrarem nos produtos fixos devem
  ser levadas a comissões estaduais.
\end{enumerate}

A não observância dessas diretrizes pode acarretar ações imediatas do
CRP~SP, sem necessidade de justificativa, consulta ou aviso prévio.
Mensagens inadequadas poderão ser encaminhadas às autoridades
responsáveis. Nosso compromisso é com a ética, a transparência e a
qualidade das nossas comunicações.

\section{Produtos}

Dividimos os produtos de comunicação do CRP~SP em \emph{produtos de
circulação rápida} e \emph{produtos de circulação lenta} (ou de ``cauda
longa''). Entre os primeiros, incluem-se conteúdo produzido para redes
sociais (como o CRP~SP Acontece e o CRP~SP na Mídia); os de circulação
lenta são aqueles com potencial para se tornar material de referência e
consulta (Cadernos Temáticos, Jornal Psi, vídeos institucionais).

A descrição de produtos de comunicação presente nesta \textbf{Política}
não impede a elaboração de diretrizes e políticas editoriais específicas
e completas para cada produto.

\subsection{Psicologia em Ação}

Divulgação de eventos a pedido, utilizando-se do \emph{site} (Agenda
CRP~SP), de Stories (perfil oficial no Instagram), postagens (demais
redes sociais) e \emph{cards} a serem compartilhados pelos territórios.
A cobertura dos eventos é publicada no CRP~SP Acontece.

As publicações de \emph{Psicologia em Ação} possuem identidade visual
própria, que pode receber elementos temáticos.

\subsection{CRP~SP Acontece}

Veiculado semanal ou quinzenalmente, informa a categoria e a sociedade
sobre as ações, eventos, representações, mobilizações, novidades etc.
promovidos nos territórios, pelas comissões gestoras, em formato de
galeria de imagens (até 10 fotos por publicação), acompanhada de breves
relatos.

Subsedes, comissões e demais instâncias devem enviar seus relatórios
preenchidos e acompanhados de foto(s) para a Comunicação.

\subsection{CRP~SP Informa}

Ferramenta de \emph{e-mail marketing} destinada a dar publicidade às
ações realizadas nos territórios, por meio de envio de mensagens
institucionais a psicólogas e entidades afins. É uma importante
ferramenta de comunicação direcionada, com média estadual de abertura
dos \emph{e-mails} de 20\%.

A cada território é facultado o envio de até quatro mensagens por mês.

\subsection{CRP~SP Articula}

Traz informações sobre parcerias do CRP~SP com outras instituições,
estatais ou de interesse público.

\subsection{CRP~SP Representa}

Divulgação da participação do CRP~SP em eventos para o qual a autarquia
foi convidada oficialmente.

\subsection{CRP~SP Debate (lives)}

Canal que possibilita levar as discussões para o público amplo e trazer
convidadas de diferentes espaços. As \emph{lives} são realizadas por
comissões e, geralmente, estão inseridas em uma campanha. Contam sempre
com tradução para Libras e têm o \emph{chat} mediado pela Coordenação de
Comunicação. São operadas no StreamYard (transmissão pelo YouTube e
Facebook), podendo apresentar até 7 participantes simultaneamente.

\emph{Lives} devem ser solicitadas com no mínimo 30 dias de
antecedência, por meio do formulário de eventos, que deve ser preenchido
e enviado para as pessoas responsáveis por eventos na Coordenação de
Comunicação. Após o recebimento, será criado grupo de imMail para
realização de testes prévios, elaboração de roteiro e alinhamento
necessário.

\subsection{Calendário de datas
comemorativas}

Artigos e postagens publicados em datas simbólicas para a defesa dos
Direitos Humanos. Textos são organizados pelas comissões permanentes e
encaminhados à ComCom com até três dias úteis de antecedência para
publicação, e devem ser elaborados com linguagem inclusiva
(\#paratodosverem), fornecendo informações sobre o significado da data
em destaque, abordando aspectos relevantes do exercício profissional da
Psicologia e fazendo referência à(s) normativa(s) do Sistema Conselhos
e/ou do Código de Ética alinhada(s) ao assunto da data. A escrita deve
ser objetiva, simples e compatível com as redes sociais (máximo de 1.800
caracteres com espaços e quatro parágrafos). O conteúdo é padronizado e
institucional.

\subsection{CRP~SP na Mídia}

Divulga as participações do CRP~SP na imprensa e em outros canais
externos.

Solicitações de divulgação devem ser encaminhadas para o \emph{e-mail}
\href{mailto:ascom@crpsp.org.br}{ascom@crpsp.org.br}}. Os pedidos
são levados para representações e plenário, para que se decida sobre o
encaminhamento.

\subsection{Comunicação integrada}

A comunicação das demais unidades administrativas do Conselho com a
categoria e com outras instituições pode ser adequada à linguagem
institucional mediante solicitação à unidade de Comunicação, observado o
cronograma e o planejamento de atividades desta unidade. A solicitação
deverá ser feita, por \emph{e-mail}, pela pessoa coordenadora da unidade
administrativa solicitante.

\subsection{Notas de pesar}

As notas de pesar têm como objetivo prestar solidariedade e
reconhecimento à pessoa falecida, especialmente quando esta desempenhou
papel significativo para a Psicologia, a defesa dos Direitos Humanos ou
outras causas alinhadas às diretrizes institucionais do CRP~SP.

A redação deve ser respeitosa, empática e compatível com a comunicação
institucional do Conselho, destacando, de forma breve, as contribuições
da pessoa homenageada. O texto não deve conter cunho político-partidário
ou religioso, garantindo imparcialidade e adequação a um público
diversificado.

\subsubsection{Solicitação e
aprovação}

As solicitações para publicação devem ser encaminhadas à Unidade de
Comunicação, acompanhadas de informações completas, como nome, data de
falecimento e breve descrição das contribuições da pessoa homenageada.

A aprovação do texto será realizada pela Comissão de Comunicação
Institucional (ComCom) ou, quando necessário, pelo Plenário.

\subsubsection{Canais de divulgação}

As notas de pesar serão divulgadas nos canais institucionais do CRP~SP,
conforme a relevância do caso e deliberação da ComCom.

Publicações de notas de pesar devem ser realizadas em caráter
excepcional, preservando a sensibilidade do público e evitando
sobrecarga nos meios de comunicação do Conselho.

\subsubsection{Prazos e responsabilidades}

A Unidade de Comunicação se compromete a publicar a nota em até
\textbf{24 horas (em dias úteis)} entre a solicitação e a validação
final.

O envio da solicitação deve ser feito com antecedência, sempre que
possível, para viabilizar o planejamento adequado.

\subsection{Informe interno}

Mensalmente, a Unidade de Comunicação publica um informe interno, com
temas relevantes para as trabalhadoras do Conselho. Textos e notícias
podem ser enviados para
\href{mailto:ascom@crpsp.org.br}{\nolinkurl{ascom@crpsp.org.br}} com o
assunto ``Informe interno'', para que seja avaliada a viabilidade de sua
publicação.

\subsection{Publicações diárias}

A Coordenação de Comunicação do CRP~SP é responsável pela publicação, a
qualquer tempo e de acordo com a demanda, de:

\begin{itemize}
\item
  material do Sistema Conselhos de Psicologia;
\item
  orientações do Centro de Referências Técnicas em Psicologia e
  Políticas Públicas (Crepop);
\item
  eventos e outros assuntos de interesse produzidos por entidades
  parceiras;
\item
  notas de posicionamento;
\item
  informes administrativos e financeiros;
\item
  orientações gerais.
\end{itemize}

\subsection{Jornal Psi}

O Jornal Psi é publicado em quatro edições anuais no formato
\emph{on-line}, com envio de versões impressas para psicólogas e
entidades com inscrição ativa no CRP~SP.

As pautas são definidas em reuniões entre a Coordenação de Comunicação e
a ComCom. Nessas ocasiões, também é definida a abordagem para cada tema.

A validação da versão final do jornal é feita pelo Plenário.

\subsubsection{Linha editorial}

Atualmente, as seções do jornal incluem:

\begin{itemize}
\item
  ``Um dia na vida'';
\item
  ``Cotidiano'';
\item
  artigo sobre tema pertinente à edição;
\item
  ``Perspectiva da/o usuária/o'';
\item
  ``Ética'';
\item
  ``Orientação'';
\item
  ``Penalidades éticas'';
\item
  ``Estante''.
\end{itemize}

\paragraph{Distribuição}

O \emph{mailing} da categoria (nomes/nomes sociais e endereços) para o
envio do Jornal Psi é gerado pelo sistema BRC, a base de dados do
Conselho. A tiragem total excede os 160 mil exemplares. Recebem o Jornal
Psi:

\begin{itemize}
\item
  psicólogas e empresas com inscrições ativas no CRP~SP;
\item
  instituições de ensino;
\item
  Sistema Conselhos de Psicologia;
\item
  bibliotecas, instituições e órgãos alinhados com a temática da
  Psicologia.
\end{itemize}

\subsection{Identidade visual da
gestão}

A Coordenação de Comunicação pode, mediante solicitação, incorporar o
mote da gestão da autarquia à identidade visual do CRP~SP.

Com o objetivo de evitar ambiguidades na comunicação, e em observância
ao princípio da impessoalidade, a partir de 2025 não produziremos
logomarcas, assinaturas gráficas ou quaisquer produtos visuais
específicos para a gestão.

\subsection{Cadernos Temáticos}

\textbf{Cadernos Temáticos} é uma série de publicações que trazem o
conteúdo de debates e eventos temáticos promovidos pelo CRP~SP. São
produzidos de acordo com a demanda das comissões permanentes e
especiais, às quais cabe a produção do conteúdo bruto e a organização
lógica do material (separação em capítulos, divisão e identificação das
falas transcritas etc.). As comissões devem garantir a qualidade e a
consistência do texto antes de encaminhar o documento à Unidade de
Comunicação.

A Unidade de Comunicação é responsável pela revisão textual, diagramação
e produção gráfica dos \textbf{Cadernos}. Não cabem à Unidade a correção
do conteúdo bruto, a revisão de transcrições e a produção de textos
originais.

\subsubsection{Atribuições}

\paragraph{Unidade de
Comunicação}

\begin{enumerate}

\item
  \textbf{Revisão gramatical e adequação de gênero}:~realizar a revisão
  gramatical e estilística do material consolidado, garantindo a
  adequação de gênero e linguagem inclusiva;
\item
  \textbf{diagramação}:~desenvolver a diagramação do conteúdo, seguindo
  as diretrizes visuais da instituição;
\item
  \textbf{produção final}:~preparar os arquivos finais para impressão ou
  publicação digital, conforme a demanda da comissão responsável;
\item
  \textbf{execução técnica}:~conduzir ajustes técnicos necessários (como
  diagramação de textos aprovados) após o recebimento do material
  validado pela comissão.
\end{enumerate}

\paragraph{Comissões solicitantes}

\begin{enumerate}

\item
  \textbf{Produção do conteúdo:}~elaborar, consolidar e revisar
  internamente o material bruto antes do envio à unidade de comunicação.
  É responsabilidade da comissão entregar o material em formato adequado
  para a revisão técnica e diagramação;
\item
  \textbf{organização e compilação:}~organizar os textos, transcrições
  (removendo inconsistências, pausas verbais e interjeições) e demais
  documentos necessários;
\item
  \textbf{validação:}~aprovar o conteúdo final revisado e diagramado
  antes da publicação ou impressão.
\end{enumerate}

\subsubsection{Estrutura}

A estrutura dos \textbf{Cadernos Temáticos} deve necessariamente conter:

\begin{enumerate}

\item
  apresentação da \textbf{coleção} Cadernos Temáticos, assinada pela
  diretoria ou pelo Plenário (até 2.500 caracteres, com espaços);
\item
  \textbf{apresentação} \textbf{da edição}, assinada pela diretoria,
  pelo Plenário ou por representante de comissão relacionada ao tema
  (até 2.000 caracteres, com espaços);
\item
  \textbf{``miolo''} --- os capítulos propriamente ditos, na forma de
  transcrições de evento, artigos, entrevistas e referências técnicas.
\end{enumerate}

Além disso, cada capítulo deve necessariamente conter:

\begin{enumerate}

\item
  título;
\item
  nome da/o autora/or ou palestrante, seguido de minicurrículo;
\item
  transcrição da fala de cada palestrante ou texto do artigo;
\item
  transcrição do debate (quando se tratar de publicação de evento), com
  identificação da pessoa que toma a palavra, sempre que possível.
\end{enumerate}

Na publicação de eventos transcritos, o mesmo título pode se referir a
diferentes falas, como acontece em mesas-redondas. Nesse caso, antes da
primeira intervenção de cada palestrante devem ser indicados nome e
minicurrículo. Nas falas subsequentes, basta indicar o nome.

Nos Cadernos Temáticos, também é comum o emprego de ``olhos'' (citações
de partes do texto em destaque).

\subsubsection{Prazo}

Uma vez recebido o texto dos \emph{Cadernos} para publicação, a
Coordenação de Comunicação deve entregar o produto finalizado (impresso
ou não) em no mínimo 30 e no máximo 60 dias.

\chapter{Serviços de comunicação e como
solicitá-los}

\section{Solicitantes}

Os serviços de comunicação e divulgação de ações, temas técnicos,
científicos e políticos nas mídias institucionais do Conselho Regional
de Psicologia de São Paulo (CRP~SP) podem ser solicitados exclusivamente
pelo Plenário, pelas comissões permanentes e especiais e pelas Unidades
Administrativas (coordenações e assessorias), em conformidade com as
diretrizes desta Política de Comunicação.

O artigo 2º, inciso I, parágrafo 1º~da
\href{https://transparencia.cfp.org.br/wp-content/uploads/sites/29/2024/06/Resolucao-03-2024.pdf}{Resolução
CRP nº~03/2024}} determina que as ações propostas por subcomissões devem
ser aprovadas junto às comissões permanentes a que se encontrem
subordinadas. \textbf{A solicitação de serviços e produtos para as
subcomissões, portanto, deve ser feita pelas comissões permanentes
responsáveis.}

\subsection{Como solicitar serviços de
comunicação}

Para requisitar serviços de comunicação, a solicitante deve:

\begin{enumerate}

\item
  preencher e encaminhar à Unidade de Comunicação o formulário
  específico para a demanda, juntamente com toda a documentação prevista
  nos respectivos formulários. É essencial atender a todos os requisitos
  indicados para que a proposta seja comprovada;
\item
  formalizar a solicitação de divulgação e/ou execução da atividade por
  meio de um processo no SEI!, dirigido à Unidade de Comunicação e com
  notificação por \emph{e-mail}. Esse setor irá prosseguir à tramitação
  e realizará as verificações cabíveis para viabilizar a solicitação
  e/ou a divulgação das atividades técnicas, científicas e políticas nos
  canais institucionais do CRP~SP.
\end{enumerate}

\section{Serviços}

\textbf{Atenção:} peças audiovisuais (\textbf{fotografias e gravações de
voz, vídeo ou ambos}) destinadas à publicação nos canais oficiais do
CRP~SP devem \textbf{obrigatoriamente} ser acompanhadas de termo de
autorização para o uso de imagem devidamente preenchido por todas as
pessoas participantes das peças, além das seguintes informações (para
cada uma das peças):

\begin{enumerate}

\item
  data e local de criação da imagem;
\item
  descrição geral do assunto (nome de evento, campanha em que se insere
  etc.);
\item
  nomes das pessoas gravadas, filmadas ou fotografadas.
\end{enumerate}

\subsection{Campanhas}

\textbf{Campanhas} são ações coordenadas que visam um objetivo comum:
por exemplo, conscientizar a categoria acerca de diretrizes éticas
emanadas do Sistema Conselhos.

As campanhas do CRP~SP têm por objetivo auxiliar na efetivação do
planejamento estratégico da autarquia. São solicitadas por comissões
permanentes e especiais e unidades administrativas.

Para toda campanha, deve ser elaborado, em parceria com a ComCom, plano
de comunicação que defina o escopo completo da campanha (mote, objetivo,
ações, linguagens, produtos, grade de veiculação etc.).

\subsubsection{Relações públicas e assessoria de
imprensa}

A Unidade de Comunicação do CRP~SP media a comunicação
interinstitucional e atua junto à imprensa, sugerindo pautas e atendendo
a solicitações de informação e entrevistas, conforme a demanda interna e
externa.

Em toda a comunicação com outras instituições e, principalmente, com a
imprensa, a Unidade de Comunicação deve definir os meios e a abordagem a
serem empregados em parceria com a ComCom.

\subsubsection{Produção e apoio a
eventos}

A Coordenação de Comunicação do CRP~SP é responsável pela gestão
(produção, cobertura e transmissão) de eventos, como mostras,
seminários, rodas de conversa, congressos, cinedebates, premiações e
eventos internos, promovidos pela autarquia \textbf{em nível estadual}.

A produção de eventos nos territórios é atribuição das respectivas
subsedes e deve ser solicitada por meio de
\href{https://forms.gle/yFMCSwpP6DQ5p9dA6}{formulário próprio}, cabendo
à Coordenação de Comunicação a orientação e a divulgação, além da
disponibilização das seguintes ferramentas digitais:

\begin{enumerate}

\item
  inclusão do evento na Agenda CRP~SP;
\item
  fornecimento de \emph{link} para inscrição;
\item
  emissão de certificados.
\end{enumerate}

\paragraph{Cobertura jornalística de
eventos}

A unidade de comunicação deve providenciar a cobertura, na forma de
fotos e vídeos, para eventos de caráter estadual. Dos eventos em
territórios, a subsede de referência é responsável pela cobertura.

Em qualquer caso, a divulgação de balanços (cobertura pós-evento) fica
condicionada ao preenchimento de
\href{https://forms.gle/aZUAWsRiPbCLaaUL7}{relatório próprio}} por
profissional qualificado para abordar o assunto de que tratou o evento e
à disponibilidade de fotografias, que devem trazer, no próprio arquivo
da imagem ou em arquivo de texto complementar:

\begin{enumerate}

\item
  data de criação da imagem;
\item
  descrição geral do assunto (nome de evento, campanha em que se insere
  etc.);
\item
  nomes das pessoas gravadas, filmadas ou fotografadas.
\end{enumerate}

\subsection{Serviços editoriais}

\subsubsection{Revisão de textos}

Os serviços editoriais prestados pela Unidade de Comunicação do CRP~SP
são:

\begin{enumerate}

\item
  revisão de texto;
\item
  preparação de originais;
\item
  criação de identidade visual e projeto gráfico;
\item
  execução de projeto gráfico;
\item
  produção gráfica/produção \emph{web}.
\end{enumerate}

O serviço de \textbf{revisão} compreende a correção de textos de acordo
com as normas gramaticais e ortográficas vigentes, além de adequação à
linguagem gendrada e demais padrões editoriais do CRP~SP. Pode, ainda,
envolver alterações no léxico, na sintaxe e na estrutura do texto, com o
objetivo de torná-lo mais acessível e/ou mais alinhado à \emph{persona}
do Conselho.

A \textbf{preparação de originais} envolve, além do serviço de revisão,
a organização de documentos mais extensos visando sua publicação na
forma de monografia (livro, apostila ou similar).

Para que os documentos produzidos possam ser publicados, é necessária a
criação de uma \textbf{identidade visual} que reflita os objetivos da
publicação e de um \textbf{projeto gráfico} subordinado a esta
identidade visual que garanta um bom padrão estético e de legibilidade.
A execução desse projeto gráfico é o que costumamos chamar de
diagramação.

A etapa seguinte à diagramação é a \textbf{produção gráfica}, e vai
desde a escolha do material e da forma de impressão à contratação de
fornecedores e à avaliação do produto final.

\subsection{Formulários e suas
funções}

\subsubsection{Formulário de solicitação e regramento de divulgação nas
mídias do
CRP~SP}

Este formulário destina-se às comissões permanentes e especiais e às
unidades administrativas que desejam solicitar divulgação institucional
nas publicações oficiais do Conselho Regional de Psicologia de São Paulo
(CRP~SP). Todas as propostas serão submetidas a uma avaliação prévia
pela Unidade de Comunicação e pela Comissão de Comunicação Institucional
(ComCom), garantindo o alinhamento com as diretrizes institucionais e a
coerência com os objetivos do CRP~SP.

O formulário deve ser utilizado para solicitar a divulgação de:

\begin{enumerate}

\item
  atividades técnicas, científicas e políticas;
\item
  divulgação de temas abordados pelo Centro de Referências Técnicas em
  Psicologia e Políticas Públicas (Crepop);
\item
  divulgação de ações de entidades parceiras;
\item
  nota de pesar;
\item
  nota de posicionamento;
\item
  informes administrativos e financeiros;
\item
  orientações.
\end{enumerate}

Esses conteúdos serão direcionados aos canais institucionais com o
objetivo de informar a categoria sobre representações e ações do
Conselho.

\subsubsection{Formulário de solicitação de
evento}

\textbf{Link:
\href{https://forms.gle/yFMCSwpP6DQ5p9dA6}{\hl{https://forms.gle/yFMCSwpP6DQ5p9dA6}}}}

O conteúdo deste formulário deve ser anexado em formato PDF ao processo
aberto no SEI, juntamente com o ``Formulário de aprovação de despesas
com atividades precípuas''. Após a aprovação no SEI, envie o número do
processo, o formulário de evento atualizado e a descrição das demandas
para a Unidade de Comunicação, setor de eventos.

\subsubsection{Relatório final do
evento}

\textbf{Link:
\href{https://forms.gle/aZUAWsRiPbCLaaUL7}{\hl{https://forms.gle/aZUAWsRiPbCLaaUL7}}}}

O Relatório Final do Evento do CRP~SP é um documento que reúne
informações essenciais sobre eventos, ações e atividades orientativas
promovidas pelo Conselho. Este relatório inclui detalhes como
identificação do evento, objetivos, atividades realizadas, avaliação dos
resultados e registros multimídia. Seu objetivo é documentar de forma
organizada e eficiente a natureza e o impacto do evento, conveniente
para fornecimento de contas, avaliações futuras e possível divulgação
nos canais do CRP~SP.

Esses procedimentos visam garantir a consistência e a efetividade das
ações de comunicação, além de garantir o alinhamento com a missão e os
valores do CRP~SP.

\chapter{Redação oficial: normas
gerais}

A redação de toda e qualquer comunicação institucional deve obedecer ao
disposto no ``\href{http://cedoc.crpsp.org.br/handle/1/2938}{Manual de
linguagem do CRP~SP}`` e normativas pertinentes. Nos casos em que não
houver definição própria do Sistema Conselhos de Psicologia, o
``\href{https://www4.planalto.gov.br/centrodeestudos/assuntos/manual-de-redacao-da-presidencia-da-republica/manual-de-redacao.pdf}{Manual
de redação da Presidência da República}`` pode ser consultado.

Para garantir a coerência na grafia de nomes, expressões e termos na
produção escrita de todos os órgãos do CFP SP, ficam instituídas as
seguintes normas gerais de redação.

\section{Nome da autarquia e de seus
órgãos}

\subsection{CRP~SP e suas subsedes}

O nome do Conselho e sua sigla podem ser grafados com referência à
divisão regional do Sistema Conselhos de Psicologia, que define São
Paulo como a 6ª região, \textbf{ou} com referência à sua área de
jurisdição (o estado de São Paulo).

Conselho Regional de Psicologia de São Paulo (CRP~SP)

\textbf{ou}

Conselho Regional de Psicologia da 6ª Região (CRP-06).

\textbf{Atenção}: a sigla deve sempre acompanhar a forma extensa.
\textbf{Não} escreva, por exemplo, ``Conselho Regional de Psicologia de
São Paulo (CRP-06)''.

\subsection{Comissões}

\begin{longtable}[]{@{}
  >{\centering\arraybackslash}p{(\linewidth - 2\tabcolsep) * \real{0.5002}}
  >{\centering\arraybackslash}p{(\linewidth - 2\tabcolsep) * \real{0.4998}}@{}}
\toprule\noalign{}
\begin{minipage}[b]{\linewidth}\centering
\textbf{Comissão}
\end{minipage} & \begin{minipage}[b]{\linewidth}\centering
\textbf{Sigla}
\end{minipage} \\
\midrule\noalign{}
\endhead
\bottomrule\noalign{}
\endlastfoot
\href{https://transparencia.cfp.org.br/crp06/legislacao/resolucao-crp-06-no-001-16-de-26-de-agosto-de-2016-cria-a-camara-de-mediacao-da-comissao-de-etica-do-conselho-regional-de-psicologia-da-6a-regiao-crp-06-cam-coe-e-aprova-seu-regulamento/}{Câmara
de Mediação da Comissão de Ética}} & CamCoe \\
\href{https://atosoficiais.com.br/cfp/resolucao-administrativa-financeira-n-14-2022-institui-e-regulamenta-o-centro-de-referencia-tecnica-em-psicologia-e-politicas-publicas-crepop-e-a-rede-crepop-como-espaco-de-operacionalizacao-das-acoes-do-crepop?origin=instituicao&q=Crepop}{Centro
de Referências Técnicas em Psicologia e Políticas Públicas}} & Crepop \\
\href{https://atosoficiais.com.br/cfp/resolucao-do-exercicio-profissional-n-3-2022-institui-condicoes-para-concessao-e-registro-de-psicologa-e-psicologo-especialistas-reconhece-as-especialidades-da-psicologia-e-revoga-as-resolucoes-cfp-no-13-de-14-de-setembro-de-2007-no-3-de-5-de-fevereiro-de-2016-e-no-8-de-25-de-abril-de-2019?origin=instituicao&q=Carpe}{Comissão
de Análise para Concessão de Registro de Psicóloga/o Especialista}} &
Carpe \\
\href{https://transparencia.cfp.org.br/crp06/legislacao/crp-sp-004-2017-cria-e-estabelece-parametros-para-o-funcionamento-da-comissao-de-auditoria-e-controle-interno-do-conselho/}{Comissão
de Auditoria e Controle Interno}} & Caci \\
\href{https://atosoficiais.com.br/cfp/resolucao-administrativa-financeira-n-5-2023-aprova-o-regimento-interno-do-conselho-regional-de-psicologia-de-sao-paulo-6a-regiao?origin=instituicao&q=05/2023}{Comissão
de Comunicação Institucional}} & ComCom \\
\href{https://atosoficiais.com.br/cfp/resolucao-administrativa-financeira-n-5-2023-aprova-o-regimento-interno-do-conselho-regional-de-psicologia-de-sao-paulo-6a-regiao?origin=instituicao&q=05/2023}{Comissão
de Ética}} & COE \\
\href{https://atosoficiais.com.br/cfp/resolucao-administrativa-financeira-n-5-2023-aprova-o-regimento-interno-do-conselho-regional-de-psicologia-de-sao-paulo-6a-regiao?origin=instituicao&q=05/2023}{Comissão
de Orientação e Fiscalização}} & COF \\
\href{https://transparencia.cfp.org.br/wp-content/uploads/sites/29/2020/08/Resolucao-CRP-002-20-Institui-a-CDHPP-final-1.pdf}{Comissão
de Direitos Humanos e Políticas Públicas}} & CDH \\
\href{https://transparencia.cfp.org.br/wp-content/uploads/sites/29/2024/06/Resolucao-03-2024.pdf}{Comissão
Especial de Entidades da Psicologia, Instituições Públicas e Coletivos
Organizados}} & Ceic \\
\href{https://transparencia.cfp.org.br/wp-content/uploads/sites/29/2020/08/Resolucao-CRP-005-20-Institui-a-Comissao-Historia-e-Memoria-final.pdf}{Comissão
de História e Memória da Psicologia}} & CHM \\
\href{https://atosoficiais.com.br/cfp/resolucao-administrativa-financeira-n-5-2023-aprova-o-regimento-interno-do-conselho-regional-de-psicologia-de-sao-paulo-6a-regiao?origin=instituicao&q=05/2023}{Comissão}
Especial} Processo Legislativo e Concursos}} & CPLC \\
\href{https://transparencia.cfp.org.br/wp-content/uploads/sites/29/2024/06/Resolucao-03-2024.pdf}{Comissão}
Especial} Psicologia Clínica e Avaliação Psicológica}} & CPAP \\
\href{https://transparencia.cfp.org.br/wp-content/uploads/sites/29/2024/06/Resolucao-03-2024.pdf}{Comissão}
Especial} Relações Étnico-Raciais}} & Crer \\
\href{https://transparencia.cfp.org.br/wp-content/uploads/sites/29/2024/06/Resolucao-03-2024.pdf}{Comissão}
Especial} Riscos, Emergências e Desastres}} & Cred \\
\href{https://atosoficiais.com.br/cfp/resolucao-administrativa-financeira-n-5-2023-aprova-o-regimento-interno-do-conselho-regional-de-psicologia-de-sao-paulo-6a-regiao?origin=instituicao&q=05/2023}{Comissão}
Permanente} de Licitação}} & CL \\
\href{https://atosoficiais.com.br/cfp/resolucao-administrativa-financeira-n-5-2023-aprova-o-regimento-interno-do-conselho-regional-de-psicologia-de-sao-paulo-6a-regiao?origin=instituicao&q=05/2023}{Ouvidoria}}
& --- \\
\end{longtable}

\subsubsection{Subcomissões}

\begin{longtable}[]{@{}
  >{\centering\arraybackslash}p{(\linewidth - 4\tabcolsep) * \real{0.3340}}
  >{\centering\arraybackslash}p{(\linewidth - 4\tabcolsep) * \real{0.3331}}
  >{\centering\arraybackslash}p{(\linewidth - 4\tabcolsep) * \real{0.3329}}@{}}
\toprule\noalign{}
\begin{minipage}[b]{\linewidth}\centering
\textbf{Comissão}
\end{minipage} & \begin{minipage}[b]{\linewidth}\centering
\textbf{Subcomissão}
\end{minipage} & \begin{minipage}[b]{\linewidth}\centering
\textbf{Sigla}
\end{minipage} \\
\midrule\noalign{}
\endhead
\bottomrule\noalign{}
\endlastfoot
COF &
\href{https://transparencia.cfp.org.br/crp06/legislacao/resolucao-crp-06-no-001-16-de-26-de-agosto-de-2016-cria-a-camara-de-mediacao-da-comissao-de-etica-do-conselho-regional-de-psicologia-da-6a-regiao-crp-06-cam-coe-e-aprova-seu-regulamento/}{Subcomissão}
Práticas Integrativas Complementares (Pics), Maconha e Psicodélicos &
--- \\
COF & Subcomissão Tics (Tecnologias de Informação e Comunicação) na
Psicologia & --- \\
COF & Subcomissão de Educação & --- \\
COF &
\href{https://transparencia.cfp.org.br/crp06/legislacao/crp-sp-004-2017-cria-e-estabelece-parametros-para-o-funcionamento-da-comissao-de-auditoria-e-controle-interno-do-conselho/}{Subcomissão}
de Psicologia do Tráfego & --- \\
CDH & Subcomissão Justiça/Prisional/Segurança Pública & --- \\
CDH & Subcomissão Drogas e Direitos Humanos & --- \\
CDH & Subcomissão Psicologia, Infâncias e Adolescências & --- \\
CDH & Subcomissão Psicologia, Pessoa com Deficiência e
Multiculturalidade & --- \\
CDH & Subcomissão Sexualidade e Gênero & --- \\
CDH & Subcomissão Mulheres & --- \\
CDH & Subcomissão Psicogerontologia da Pessoa Idosa & --- \\
\end{longtable}

\subsection{Unidades administrativas}

\begin{longtable}[]{@{}
  >{\centering\arraybackslash}p{(\linewidth - 2\tabcolsep) * \real{0.5002}}
  >{\centering\arraybackslash}p{(\linewidth - 2\tabcolsep) * \real{0.4998}}@{}}
\toprule\noalign{}
\begin{minipage}[b]{\linewidth}\centering
\textbf{Unidade}
\end{minipage} & \begin{minipage}[b]{\linewidth}\centering
\textbf{Sigla}
\end{minipage} \\
\midrule\noalign{}
\endhead
\bottomrule\noalign{}
\endlastfoot
Jurídico & --- \\
Recursos Humanos & RH \\
Secretaria & --- \\
\end{longtable}

\subsubsection{Gerências}

\begin{longtable}[]{@{}
  >{\centering\arraybackslash}p{(\linewidth - 2\tabcolsep) * \real{0.5002}}
  >{\centering\arraybackslash}p{(\linewidth - 2\tabcolsep) * \real{0.4998}}@{}}
\toprule\noalign{}
\begin{minipage}[b]{\linewidth}\centering
\textbf{Gerência}
\end{minipage} & \begin{minipage}[b]{\linewidth}\centering
\textbf{Sigla}
\end{minipage} \\
\midrule\noalign{}
\endhead
\bottomrule\noalign{}
\endlastfoot
Gerência de Administração e Tecnologia da Informação & Gati \\
Gerência Técnica e Política & GTP \\
Gerência de Relações Institucionais & GRI \\
\end{longtable}

\subsubsection{Coordenações}

\begin{longtable}[]{@{}
  >{\centering\arraybackslash}p{(\linewidth - 2\tabcolsep) * \real{0.5002}}
  >{\centering\arraybackslash}p{(\linewidth - 2\tabcolsep) * \real{0.4998}}@{}}
\toprule\noalign{}
\begin{minipage}[b]{\linewidth}\centering
\textbf{Coordenação}
\end{minipage} & \begin{minipage}[b]{\linewidth}\centering
\textbf{Sigla}
\end{minipage} \\
\midrule\noalign{}
\endhead
\bottomrule\noalign{}
\endlastfoot
Coordenação de Gestão Administrativa & CGA \\
Coordenação de Tecnologia da Informação & CTI \\
Coordenação de Gestão Financeira & CGF \\
Coordenação de Comunicação & CCOM \\
Coordenação de Atendimento & CA \\
Coordenação de Ética & CE \\
Coordenação de Orientação e Fiscalização & \\
Coordenação de Subsedes & CS \\
\end{longtable}

\chapter{Apêndice A --- Diretrizes para a comunicação de conselheiras e
trabalhadoras no CRP~SP nas redes
sociais}

A presença do Conselho Regional de Psicologia de São Paulo (CRP~SP) nas
redes sociais dá-se no cumprimento das atribuições da autarquia de
orientar e informar a categoria sobre temas associados à Psicologia, à
atuação profissional e ao funcionamento do Sistema Conselhos de
Psicologia. Objetiva complementar os canais proprietários da
instituição, buscando ampliar seus pontos de contato e comunicação com
psicólogas de todo o estado de São Paulo, assim como participar dos
debates, em sociedade e com demais entidades, sobre a Psicologia como
ciência e profissão.

As conselheiras em vigência de mandato representam o CRP~SP perante a
sociedade a qualquer tempo. Por isso, suas manifestações públicas ---
palestras, entrevistas a veículos de imprensa, comentários em redes
sociais etc. --- devem ser pautadas pelos princípios dispostos nesta
\textbf{Política}.

Fica definido que:

\begin{enumerate}

\item
  só serão permitidos perfis e páginas associados ao CRP~SP que sejam
  criados e mantidos pela instituição, ou seja, tenham caráter de
  oficiais. A autarquia reserva-se o direito de não corroborar perfis e
  páginas que se autointitulem associados ao Conselho e não sejam
  reconhecidos como perfis oficiais, tomando as devidas medidas quando
  necessárias;
\item
  as contas oficiais do CRP~SP serão gerenciadas por equipe certificada
  e autorizada pelo CRP~SP, composta por profissionais da comunicação,
  trabalhadoras e trabalhadores da autarquia e psicólogas e psicólogos
  conselheiras e conselheiros;
\item
  como forma de garantia da veracidade, exatidão, confiabilidade,
  qualidade e legalidade das informações e mensagens que circulam pelos
  canais e perfis oficiais do CRP~SP, as interações, nestes espaços,
  estarão sujeitas à moderação e filtragem, podendo ser excluídas as
  mensagens e os perfis que

  \begin{enumerate}
  
  \item
    usem linguagem inapropriada, obscena, caluniosa, grosseira, abusiva,
    difamatória, ofensiva ou de qualquer outra forma censurável;
  \item
    façam apologia a práticas ilícitas;
  \item
    incitem o ódio, a violência, o racismo ou façam discriminação de
    qualquer ordem;
  \item
    contenham ameaças, assédio, injúria, calúnia ou difamação ou
    configurem qualquer outra forma de ilícito penal;
  \item
    divulguem conteúdos na forma de \emph{spam} ou ``correntes'';
  \item
    caracterizem intuito comercial ou publicitário;
  \item
    sejam ininteligíveis ou fora de contexto;
  \item
    contenham propagandas político-partidárias;
  \item
    contenham \emph{links} suspeitos ou representem ameaça à segurança
    da informação;
  \item
    façam uso de informações e imagens de pessoas e instituições de modo
    indevido;
  \item
    contenham dados pessoais da autora, do autor ou de terceiros;
  \item
    violem os direitos de imagem e de propriedade intelectual;
  \item
    sejam fraudulentos (perfis falsos e/ou \emph{bots}) ou promovam
    conteúdo inverídico (\emph{fake news}).
  \end{enumerate}
\end{enumerate}

A não observância a essas regras poderá acarretar ações imediatas do
CRP~SP, independentemente de justificativa, consulta ou aviso prévio.
Conforme o conteúdo, as mensagens poderão ser encaminhadas à autoridade
responsável.

\section{Orientações gerais}

\textbf{O ciberespaço é um espaço público}. Todas as opiniões,
comentários e interações emitidas/ocorridas em redes sociais, portanto,
devem ser entendidas como manifestações públicas, sujeitas à
responsabilização na forma da lei.

A participação do Sistema Conselhos de Psicologia na arena pública,
tanto em sua dimensão virtual quanto física, é estritamente
institucional. Por isso, recomenda-se o máximo de cuidado a conselheiras
e trabalhadoras do CRP~SP antes de emitir qualquer posicionamento,
comentário ou opinião pessoal sobre assuntos que digam respeito ao
Sistema Conselhos e/ou à prática da Psicologia. Além disso,
\textbf{perfis pessoais são expressamente vedados de falar em nome do
CRP~SP}. Menções à autarquia, a conselheiras e trabalhadoras, desde que
não violem questões de direitos autorais, privacidade e segurança da
informação, e que sejam feitas com respeito e responsabilidade, são
autorizadas.

Caso seja necessário emitir um posicionamento e/ou resposta oficial
sobre qualquer assunto, a Unidade de Comunicação deve ser acionada por
um solicitante legítimo. É fundamental que o conteúdo enviado esteja
previamente validado por uma comissão permanente ou especial ou por uma
unidade administrativa, garantindo que o tema seja tratado de forma
\textbf{técnica} e alinhada aos parâmetros institucionais.

Toda e qualquer comunicação pessoal deve ser tratada como tal. Mensagens
de caráter pessoal não devem ser tornadas públicas. \textbf{Havendo
necessidade de divulgar o conteúdo de alguma mensagem de terceiros,
certifique-se de obter uma autorização prévia da autora}. Tendo em vista
a possibilidade de divulgação, mantenha, em sua correspondência
institucional, o mesmo respeito e urbanidade que caracterizam as
manifestações públicas do CRP~SP.

\textbf{No engajamento com postagens não relacionadas ao Sistema
Conselhos e/ou à prática da Psicologia, conselheiras e trabalhadoras
devem se portar de maneira respeitosa e condizente com os princípios
definidos nas normativas do Sistema Conselhos e, mais amplamente, com o
ordenamento jurídico brasileiro}. Ressaltamos a importância de não
disseminar nenhum conteúdo que não proceda de fonte identificável,
confiável e verificável.

\section{\texorpdfstring{Dicas de comunicação
\emph{on-line}}{Dicas de comunicação on-line}}

A seguir, listamos uma série de dicas adaptadas do
\href{https://www.gov.br/economia/pt-br/centrais-de-conteudo/publicacoes/guias-e-manuais/ManualdeRedesSociais.pdf}{Manual
de redes sociais}} do Ministério da Economia, para que sejam usadas como
referência.

Os conteúdos postados são sempre de ordem pessoal, mas, a partir do
momento em que a usuária definir seu local de trabalho, invariavelmente
terão também um teor institucional/profissional. Isso significa que há
algumas boas práticas simples que devem ser seguidas quando a
conselheira ou trabalhadora postar alguma imagem ou texto de seu local
de trabalho ou com referência a ele:

\begin{itemize}
\item
  as informações escritas na rede são de responsabilidade da
  conselheira/trabalhadora, mas atingem a todas as entidades --- pessoas
  e instituições --- presentes nos conteúdos de suas mensagens;
\item
  você é uma pessoa pública. Sempre que postar algo nas redes sociais,
  entenda que o conteúdo da sua mensagem será visto por colegas, chefes,
  clientes, fornecedores, parceiras institucionais e de negócio, amigas
  e familiares;
\item
  suas seguidoras/amigas vão confundir o seu ``eu'' pessoal com o seu
  ``eu'' profissional. Você pode não ser a porta-voz oficial da
  instituição na qual trabalha, mas, a partir do momento em que deixar
  público seu perfil numa rede social, será vista pelas demais usuários
  como alguém que fala em nome da instituição;
\item
  evite postar qualquer coisa que possa gerar danos à instituição em que
  atue;
\item
  escrever na rede é o mesmo que escrever em pedra. Escrever não é o
  mesmo que falar: suas palavras ficam na \emph{web} e são indexadas
  quase que instantaneamente por outras redes. Assim, mesmo que apague
  um \emph{post} do qual tenha se arrependido, ele provavelmente já terá
  sido indexado pelo Google e por outros \emph{sites}, eternizando-se na
  internet e pondo-se ao alcance de todos os usuários. Pense antes de
  publicar. Se for para se arrepender, arrependa-se antes de escrever;
\item
  proteja-se. Uma crise envolvendo \emph{posts} em redes sociais feitos
  por uma trabalhadora e prejudicando o órgão nunca tem somente a
  instituição como alvo único, e pode acabar causando danos de imagem
  também a suas funcionárias.
\item
  nunca deixe de ser você. Como qualquer cidadã, você é livre para
  pensar e expressar o que desejar, da forma que preferir. Mas, como
  qualquer pessoa pública, tem de entender que tudo o que expressar
  provavelmente trará consequências, sejam positivas ou negativas;
\item
  cuidado com questões sigilosas e informações privilegiadas. É vedada a
  publicação de informação confidencial ou prejudicial à imagem do
  Conselho, assim como não se deve antecipar resultados de trabalhos
  ainda não divulgados oficialmente pela autarquia;
\item
  tenha uma conduta ética nas redes sociais durante viagens de trabalho.
  Conselheiras e/ou trabalhadoras que estiverem em viagem a serviço do
  Conselho devem tomar cuidado com suas postagens em redes sociais. Há
  diversos casos em que a imagem de órgãos públicos foi arranhada devido
  a mau comportamento de servidores em viagens oficiais que é exposto
  nas redes;
\item
  lembre-se: você não está a turismo. Ainda que nas horas vagas você
  aproveite uma ida pelo CRP~SP ao Rio de Janeiro, por exemplo, e
  resolva pegar uma praia ou visitar um ponto turístico, mostrar isso
  nas redes pode gerar um conflito de mensagens para seus seguidores e
  principalmente para colegas de trabalho e amigos que sabem do real
  propósito de sua missão. As viagens são custeadas com recursos
  públicos e isso deve ser levado em máxima consideração;
\item
  mostre a que veio. Postagens sobre o trabalho realizado, de forma
  positiva, são salutares e até incentivadas. Divulgar o nosso trabalho
  e o trabalho do CRP~SP é sempre positivo e colabora para a boa imagem
  da instituição.
\end{itemize}

\chapter{Apêndice B --- Modelo de autorização para uso de imagem, voz e
textos em publicações impressas e
eletrônicas}

Eu,
\_\_\_\_\_\_\_\_\_\_\_\_\_\_\_\_\_\_\_\_\_\_\_\_\_\_\_\_\_\_\_\_\_\_\_\_\_\_\_\_\_\_\_\_\_\_\_\_\_\_\_\_\_\_\_\_\_\_\_\_\_\_\_\_\_\_\_\_\_\_\_\_,
portadora/or do RG
nº~\_\_\_\_\_\_\_\_\_\_\_\_\_\_\_\_\_\_\_\_\_\_\_\_\_\_\_, residente e
domiciliada/o à
\_\_\_\_\_\_\_\_\_\_\_\_\_\_\_\_\_\_\_\_\_\_\_\_\_\_\_\_\_\_\_\_\_\_\_\_\_\_\_\_\_\_\_\_\_\_\_\_\_\_\_\_\_\_\_\_\_\_\_\_\_\_\_\_\_\_\_\_\_\_\_\_\_\_\_,
telefone ( \_\_\_ ) \_\_\_\_\_\_\_\_\_\_\_\_\_\_\_\_\_\_\_\_\_\_\_,
celular ( \_\_\_ ) \_\_\_\_\_\_\_\_\_\_\_\_\_\_\_\_\_\_\_\_\_\_\_\_\_,
\emph{e-mail}
\_\_\_\_\_\_\_\_\_\_\_\_\_\_\_\_\_\_\_\_\_\_\_\_\_\_\_\_\_\_\_\_\_\_\_\_\_\_\_\_,
autorizo o Conselho Regional de Psicologia da 6ª Região (CRP 06), com
sede à Rua Teodoro da Silva, 417, Jd. América, São Paulo (SP), a
produzir, publicar e distribuir minha imagem e voz em todos os meios de
comunicação e/ou publicações impressas, quando de minha participação em
entrevista, palestra/apresentação e/ou comentários realizados em evento
no dia \_\_\_ / \_\_\_ / \_\_\_\_\_.

Declaro estar ciente de que a referida autorização é concedida de forma
gratuita e por prazo indeterminado, abrangendo todas as mídias e
formatos existentes ou que venham a ser criados, podendo ser utilizada
em campanhas publicitárias, materiais promocionais, redes sociais e/ou
\emph{website} do CRP~SP, entre outros meios de divulgação.

Por meio desta autorização, renuncio expressamente a qualquer
remuneração, pagamento, indenização ou compensação financeira pelo uso
da minha imagem e voz, concordando em ceder os direitos de uso de forma
gratuita e irrestrita.

Declaro ainda estar ciente de que a presente autorização é irrevogável e
intransferível, sendo válida mesmo após o término de meu vínculo com o
CRP~SP.

Por fim, declaro estar ciente dos termos desta autorização e concordar
com todas as disposições nela contidas.

Cidade, dia, mês e ano:

\_\_\_\_\_\_\_\_\_\_\_\_\_\_\_\_\_\_\_\_\_\_\_\_\_\_\_\_\_\_\_\_\_\_\_,
\_\_\_\_ de \_\_\_\_\_\_\_\_\_\_\_\_\_\_\_\_\_\_ de \_\_\_\_\_\_\_.

\_\_\_\_\_\_\_\_\_\_\_\_\_\_\_\_\_\_\_\_\_\_\_\_\_\_\_\_\_\_\_\_\_

\textbf{Assinatura}

\end{document}
